\documentclass[12pt]{article}
\usepackage{amsmath, amssymb}
\usepackage[margin=1in]{geometry}
\setlength{\parskip}{6pt}
\title{QUBO Formulation: Quantum-Circuit Qubit Allocation}
\date{}
\begin{document}
\maketitle

\subsection*{Quantum-Circuit Qubit Allocation}

\paragraph{Problem Statement.}
Given a set of logical qubits $L$, a set of physical qubits $P$, and a hardware coupling graph that defines the connectivity and fidelity between physical qubits, the task is to assign each logical qubit to a distinct physical qubit. Let $G \subseteq L \times L$ denote the set of logical qubit pairs that participate in two-qubit gates, and let $C_{p,q}$ denote the penalty cost associated with mapping a logical pair to physical qubits $p$ and $q$. The goal is to find a one-to-one mapping that minimizes the total penalty incurred by two-qubit gates mapped onto distant or low-fidelity physical pairs.

\paragraph{Decision Variables.}
For each logical qubit $l \in L$ and physical qubit $p \in P$, we define a binary decision variable $x_{l,p} \in \{0,1\}$. The variable $x_{l,p}$ equals $1$ if logical qubit $l$ is mapped to physical qubit $p$, and $0$ otherwise. The total number of variables is $|L| \times |P|$.

\paragraph{QUBO Derivation.}
\textbf{Step 1 --- Objective Function.}
The objective is to minimize the total penalty for two-qubit gates mapped to physical pairs. This is expressed as:
\begin{align}
\min_{x} \quad J(x) = \sum_{(l_i, l_j) \in G} \sum_{p \in P} \sum_{q \in P} C_{p,q} x_{l_i,p} x_{l_j,q}.
\end{align}

\textbf{Step 2 --- Constraints.}
The mapping must satisfy two fundamental constraints. First, each logical qubit must be assigned to exactly one physical qubit:
\begin{align}
\sum_{p \in P} x_{l,p} = 1, \quad \forall l \in L.
\end{align}
Second, each physical qubit can host at most one logical qubit:
\begin{align}
\sum_{l \in L} x_{l,p} \le 1, \quad \forall p \in P.
\end{align}

\textbf{Step 3 --- Penalty Term Conversion.}
We convert the constraints into quadratic penalty terms added to the objective. Let $\lambda > 0$ be a penalty weight.

(a) For the logical assignment constraint, we introduce:
\begin{align}
H_{\text{C1}}(x) = \lambda \sum_{l \in L} \left( \sum_{p \in P} x_{l,p} - 1 \right)^2.
\end{align}
(b) Expanding the square and applying the idempotent property $x_{l,p}^2 = x_{l,p}$ for binary variables yields:
\begin{align}
\left( \sum_{p \in P} x_{l,p} - 1 \right)^2 &= \sum_{p \in P} x_{l,p}^2 - 2 \sum_{p \in P} x_{l,p} + 1 \nonumber \\
&= \sum_{p \in P} x_{l,p} - 2 \sum_{p \in P} x_{l,p} + 1 \nonumber \\
&= 1 - \sum_{p \in P} x_{l,p}.
\end{align}
(c) The resulting general penalty contribution is:
\begin{align}
H_{\text{C1}}(x) = \lambda \sum_{l \in L} \left( 1 - \sum_{p \in P} x_{l,p} \right) = \lambda n_L - \lambda \sum_{l \in L} \sum_{p \in P} x_{l,p}.
\end{align}

For the physical capacity constraint, we penalize any assignment where two or more logical qubits occupy the same physical qubit:
\begin{align}
H_{\text{C2}}(x) = \lambda \sum_{p \in P} \left( \sum_{l \in L} x_{l,p} \right)^2.
\end{align}
Expanding and applying idempotence:
\begin{align}
\left( \sum_{l \in L} x_{l,p} \right)^2 &= \sum_{l \in L} x_{l,p}^2 + \sum_{l \in L} \sum_{l' \in L, l' \neq l} x_{l,p} x_{l',p} \nonumber \\
&= \sum_{l \in L} x_{l,p} + \sum_{l \neq l' \in L} x_{l,p} x_{l',p}.
\end{align}
The resulting penalty contribution is:
\begin{align}
H_{\text{C2}}(x) = \lambda \sum_{p \in P} \left( \sum_{l \in L} x_{l,p} + \sum_{l \neq l' \in L} x_{l,p} x_{l',p} \right).
\end{align}

\textbf{Step 4 --- Combined QUBO Hamiltonian.}
Summing the objective and penalty terms gives the total Hamiltonian $H(x)$:
\begin{align}
H(x) = \sum_{(l_i, l_j) \in G} \sum_{p \in P} \sum_{q \in P} C_{p,q} x_{l_i,p} x_{l_j,q} + \lambda \sum_{l \in L} \sum_{p \in P} x_{l,p} + \lambda \sum_{p \in P} \sum_{l \neq l' \in L} x_{l,p} x_{l',p} + \lambda n_L.
\end{align}
Note that the linear terms from $H_{\text{C1}}$ and $H_{\text{C2}}$ cancel exactly, leaving only quadratic and constant terms.

\textbf{Step 5 --- Q Matrix Entries and Offset.}
Reading off the coefficients from $H(x) = \sum_{i \le j} Q_{ij} x_i x_j + c$, we obtain the following closed-form expressions for the QUBO matrix entries and offset:
\begin{align}
Q_{(l_i,p), (l_j,q)} &= C_{p,q} & \text{if } (l_i, l_j) \in G \text{ and } (l_i,p) \neq (l_j,q), \\
Q_{(l,p), (l',p)} &= \lambda & \text{if } l \neq l', \\
Q_{(l,p), (l,p)} &= 0 & \text{for all } l \in L, p \in P, \\
c &= \lambda n_L.
\end{align}

\paragraph{Penalty Weight Justification.}
The penalty weight $\lambda$ must be chosen such that any constraint violation incurs a cost strictly greater than the maximum possible reduction in the objective function. The objective function is bounded by $|G| \cdot \max_{p,q} |C_{p,q}|$. Therefore, setting $\lambda > |G| \cdot \max_{p,q} |C_{p,q}|$ guarantees that the energy of any infeasible configuration exceeds that of the optimal feasible configuration, ensuring that the optimizer prioritizes feasibility over objective minimization.

\paragraph{Correctness Argument.}
Minimizing $H(x)$ over the binary hypercube $\{0,1\}^{|L| \times |P|}$ recovers the optimal qubit allocation because feasible mappings satisfy all constraints, yielding zero penalty terms and reducing $H(x)$ to the original objective $J(x)$. Conversely, any infeasible assignment violates at least one constraint, incurring a penalty strictly larger than the maximum possible objective gain from that violation. Consequently, the global minimum of $H(x)$ must correspond to a feasible assignment with the minimum allocation cost.

\paragraph{Illustrative Example.}
Consider a small instance with $n_L = 2$ logical qubits ($L = \{0,1\}$) and $n_P = 2$ physical qubits ($P = \{0,1\}$). Suppose there is a single two-qubit gate between logical qubits $0$ and $1$, so $G = \{(0,1)\}$. The variable mapping is $x_{0,0} \to x_0$, $x_{1,0} \to x_1$, $x_{0,1} \to x_2$, $x_{1,1} \to x_3$. Using the general formulas with penalty weight $\lambda = 100$ and offset $c = 200$, the objective terms are:
\begin{align}
H_{\text{obj}} = C_{0,0} x_0 x_1 + C_{0,1} x_0 x_3 + C_{1,0} x_1 x_2 + C_{1,1} x_2 x_3.
\end{align}
The penalty terms expand to:
\begin{align}
H_{\text{C1}} &= 100(1 - x_0 - x_2) + 100(1 - x_1 - x_3) = 200 - 100(x_0 + x_1 + x_2 + x_3), \\
H_{\text{C2}} &= 100(x_0 x_2 + x_1 x_3) + 100(x_0 + x_2) + 100(x_1 + x_3).
\end{align}
Combining these, the linear terms cancel, yielding the full Hamiltonian:
\begin{align}
H(x) = C_{0,0} x_0 x_1 + C_{0,1} x_0 x_3 + C_{1,0} x_1 x_2 + C_{1,1} x_2 x_3 + 100 x_0 x_2 + 100 x_1 x_3 + 200.
\end{align}
The optimal mapping assigns logical qubit $0$ to physical qubit $0$ and logical qubit $1$ to physical qubit $1$ (i.e., $x_0=1, x_3=1$, others $0$), yielding an energy of $C_{0,1} + 200$. Alternatively, mapping $0 \to 1$ and $1 \to 0$ gives $C_{1,0} + 200$. The minimum energy corresponds to the physical pair with the lowest gate penalty, demonstrating how the QUBO formulation correctly encodes the allocation task.

\end{document}